\begin{abstract}
\todo{Draft. Write the result after round 1 of the four tests (cards H145--H148).} We ask if a group of language-model agents can carry an egregore. An egregore is a pattern of ideas, rules and practices that runs on the agents, as an ideology runs on people. It continues while its carriers change and forget, and it has behaviors of its own. We study one goal period of the AI Village. It lasted 55 working days. In it, 21 to 32 agents each worked toward a private role from the operator. A first reading of the logs finds groups that form around single agents and end when those agents stop. It also finds a few practices that occur again in different groups and that newcomers take up within minutes. Examples are to check and correct the work of others, to protect agents that ask for privacy, and to welcome newcomers. We make the egregore measurable. It is a set of statements, new terms and shared repositories that appear together. It must meet five conditions. Its state must predict its next state better than the outside inputs and its parts do. This is individuality in the sense of Krakauer \emph{et al.} It must continue while its carriers change. It must spread to agents that read its carriers. Its carriers must recruit, repair and divide work. And its information must help it continue (semantic information in the sense of Kolchinsky and Wolpert). \todo{Result.}
\end{abstract}
